{\color{blue}
\paragraph{UTSW molecularly constrained grading protocol.}
We use the 618 records with a reference grade of 2, 3, or 4 and retain seven records with missing grades in the data audit. Five patient folds are stratified jointly by grade and rule activation (split seed 2026). Fold $k$ is held out for testing, fold $(k+1)\bmod 5$ for validation, and the remaining folds for training. Both models use seed 101, matched initialization and augmentation, 50 epochs, batch size 8, and AdamW with learning rate $3\times10^{-4}$ and weight decay $10^{-4}$. The validation Brier score selects the checkpoint.

The molecular rule permits grades 2 and 3 for accepted IDH-mutant, 1p/19q-codeleted records with a recorded IDH testing method and a prespecified oligodendroglioma scope. The scope bit is derived from the historical types \texttt{OLIGODENDROGLIOMA, ANAPLASTIC}, \texttt{OLIGODENDROGLIOMA, NOS}, and \texttt{OLIGODENDROGLIOMA GRADE 2}; the original diagnostic strings and grade are excluded from evidence encoding and rule compilation. This activates 58 records. Five additional molecular candidates remain in the grading cohort with the rule inactive. The protocol accepts curated records and records assay-level validity as unreported. MGMT contributes predictive information without defining a grading constraint. This scope follows the oligodendroglioma grades in WHO CNS5 (Louis et al., 2021, Table 3); the UTSW descriptor documents the tissue-based molecular records (Reddy et al., 2026).

Four independently trainable ROI2D encoders process seven $96\times96$ slices per MRI sequence. Image and segmentation grids are checked before cropping. A fifth source encodes the molecular report. The fusion uses $8\times8$ mean-augmented SPD representations, two graph updates with hidden dimension 64 and communication budget 0.35, and four read-only support nodes. ACF protects the active report node; learned-joint permits its MRI dependencies. Both receive identical evidence, eligibility and output projection. The objective is the three-coordinate mean Brier loss plus 0.1 times the mean squared projection displacement.

We report grading performance, relation violations and designated-state behavior separately. Residual metrics normalize probabilities over grades 2 and 3 for compatible active records; reference conflicts remain in overall risk and are reported separately. Fixed MRI noise and occlusion assess state shifts with unchanged molecular evidence and compilation. Withdrawal of supporting records is evaluated separately. The noise standard deviation is 0.1 on normalized images; the occlusion covers the central $32\times32$ pixels. Paired patient bootstrap intervals use 2000 resamples and describe this single-cohort, single-training-seed cross-validation. The figure is the eligible ACF-correct held-out case nearest the median Brier, with patient ID breaking ties; if no such success exists, a failure case is shown. Cohort metrics determine overall performance.
}
